\scanpage{11}
5. Soit tjrs $M^\bullet\in\operatorname{Ob}D^+(X/S)$, et considérons le complexe $Q^0[M^\bullet]$ \deleted{de} \deleted{\unsure{de $Q^0[$}} de (pro-$\underline O_X$-Modules) à stratification, dont la connaissance implique celle de $M^\bullet$ : du pt de vue de \unsure{complexe} additif d'abord, en passant aux faisceaux des sections stationnaires, puis les $\underline O_X$-structures sur les $M^i$ comme restrictions des structures extrême-droites des $Q^0[M^i]$ (lesquelles, rappelons le encore, ne sont cependant pas respectées par l'op. diff. de $Q^0[M^\bullet]$).

Considérons les $\mathcal H^i(\ell Q^0[M^\bullet])$, et supposons que leur formation commute aux changts de base \deleted{\unsure{semi}}\inserted{co}-simpliciaux $\mathcal P^0\longrightarrow\mathcal P^{\deleted{\unsure{1}}n}$ ; ceci sera assuré par des hypothèses
\margin{10}
